\uuid{nZ3g}
\exo7id{4141}
\titre{exo7 4141}
\auteur{quercia}
\organisation{exo7}
\datecreate{2010-03-11}
\isIndication{false}
\isCorrection{true}
\chapitre{Fonction de plusieurs variables}
\sousChapitre{Dérivée partielle}
\cpge{mpsi-mp2i.fonctions-deux-variables.b, pcsi.fonctions-deux-variables.b, ptsi.fonctions-deux-variables.b, mpsi-mp2i.fonctions-deux-variables.c, pcsi.fonctions-deux-variables.c, ptsi.fonctions-deux-variables.c, bcpst1.fonctions-deux-variables.c, bcpst1.fonctions-deux-variables.d, mp-mpi.calcul-differentiel-optimisation.a}
\module{Analyse}
\niveau{L2}
\difficulte{}

\contenu{
\question{
Soient ${u,v,f,g} : {\R^2} \mapsto \R$ des fonctions de classe $\mathcal{C}^2$ liées
par la relation :
$$\forall\ (x,y) \in {\R^2},\ f(x,y) = g(u(x,y),v(x,y)).$$

Calculer les dérivées partielles premières et secondes de $f$
en fonction de celles de $g$.
}
\reponse{
$\frac{\partial^2 f}{\partial x \partial y}= \frac{\partial u}{\partial x}\frac{\partial u}{\partial y}\frac{\partial^2 g}{\partial u^2} + \left(\frac{\partial u}{\partial x}\frac{\partial v}{\partial y}+\frac{\partial u}{\partial y}\frac{\partial v}{\partial x}\right)\frac{\partial^2 g}{\partial u \partial v}
         + \frac{\partial v}{\partial x}\frac{\partial v}{\partial y}\frac{\partial^2 g}{\partial v^2} + \frac{\partial^2 u}{\partial x \partial y}\frac{\partial g}{\partial u} + \frac{\partial^2 u}{\partial x \partial y}\frac{\partial g}{\partial v}$.
}
}
